\documentclass[a4paper,14pt]{extarticle}
\usepackage[utf8]{inputenc}
\usepackage[russian]{babel}
\usepackage[T1,T2A]{fontenc}
\usepackage{geometry}
\geometry{left=30mm, right=15mm, top=20mm, bottom=20mm}
\usepackage{hyperref}
\usepackage{enumitem}
\usepackage{titlesec}
\usepackage{setspace}
\onehalfspacing
\usepackage{indentfirst}

\begin{document}

\begin{center}
    \textbf{\Large Оптимизация поиска базовых чанков в подходах SBC}
\end{center}

\vspace{1em}

\section*{Аннотация}
В данной работе рассматривается проблема оптимизации поиска базовых чанков в системах дедупликации данных на основе сходства (Similarity-Based Chunking, SBC). Традиционные подходы к дедупликации (например, Content-Defined Chunking) позволяют устранять лишь полные дубликаты блоков данных, тогда как методы SBC расширяют этот подход, находя похожие (но не идентичные) блоки и применяя к ним дельта-компрессию. Это позволяет значительно повысить итоговый коэффициент сжатия.

Однако поиск наиболее подходящих базовых чанков для дельта-сжатия среди огромного массива данных является вычислительно сложной задачей. Применение полного побайтового сравнения нецелесообразно из-за высоких накладных расходов на процессор и память. Существующие эвристики, такие как хеширование Ароновича, алгоритмы Odess или Broder's method, предлагают компромисс между скоростью поиска и степенью сжатия за счёт использования контентно-зависимой выборки (Content-Defined Sampling) и построения легковесных эскизов (sketches).

Предполагаемые результаты работы включают разработку новых или улучшение существующих эвристических алгоритмов в рамках системы \texttt{marline} (реализованной на языке Rust), а также проведение сравнительного анализа их производительности и степени дедупликации на репрезентативных наборах данных. Ожидается, что предложенные оптимизации позволят снизить вычислительную нагрузку при поиске базовых чанков, сохраняя или улучшая итоговый коэффициент дедупликации.

\newpage

\section*{Формальная цель и декомпозиция на задачи}

\textbf{Формальная цель:} Исследование и разработка оптимизированных эвристических алгоритмов поиска базовых чанков в системах дедупликации данных на основе сходства (SBC) для повышения производительности системы без существенной потери коэффициента сжатия.

\textbf{Декомпозиция на задачи:}

\begin{enumerate}
    \item \textbf{Анализ предметной области и обзор литературы}
    \begin{itemize}
        \item Изучение базовых алгоритмов дедупликации (CDC, SBC).
        \item Анализ существующих методов поиска сходства (хеширование Ароновича, Odess, N-Transform, методы на основе графов).
        \item Подбор и структурирование литературы.
    \end{itemize}
    
    \item \textbf{Изучение кодовой базы проекта \texttt{marline}}
    \begin{itemize}
        \item Архитектурный анализ подсистем \texttt{marline\_scrub} и \texttt{marline\_sketcher}.
        \item Профилирование текущих реализаций эвристик для выявления вычислительных узких мест при выборе базовых чанков.
    \end{itemize}
    
    \item \textbf{Разработка и внедрение эвристик поиска базовых чанков}
    \begin{itemize}
        \item Проектирование новых эвристик (включая фильтрацию кандидатов и улучшенное сэмплирование).
        \item Программная реализация алгоритмов на языке Rust.
        \item Интеграция эвристик в существующий конвейер дедупликации.
    \end{itemize}
    
    \item \textbf{Тестирование и бенчмаркинг}
    \begin{itemize}
        \item Подготовка датасетов (репрезентативных наборов данных).
        \item Проведение сравнительных тестов: замер коэффициента дедупликации (CDC vs SBC), времени выполнения, потребления оперативной памяти.
    \end{itemize}
    
    \item \textbf{Анализ результатов и подготовка отчётности}
    \begin{itemize}
        \item Систематизация полученных метрик и оценка эффективности внедрённых оптимизаций.
        \item Оформление итогового отчёта и презентации по стандартам ГОСТ.
    \end{itemize}
\end{enumerate}

\newpage

\section*{Обзор литературы}

\subsection*{1. Lior Aronovich et al. "The design of a similarity based deduplication system" (SYSTOR '09)}
\textbf{Ссылка:} \url{https://dl.acm.org/doi/10.1145/1534530.1534541} \\
\textbf{Нейросетевой обзор:} В статье представлена архитектура одной из первых промышленных систем дедупликации на основе сходства. Авторы предлагают метод вычисления специальных признаков (features) для каждого блока данных, что позволяет находить похожие чанки, а не только точные копии, после чего к ним применяется дельта-кодирование. \\
\textbf{Зачем это нужно в работе:} Это базовая (основополагающая) статья для понимания концепции SBC и хеширования Ароновича, которое уже частично реализовано в проекте \texttt{marline}. Помогает задать теоретический фундамент и определить baseline для сравнения эвристик.

\subsection*{2. "Odess: A Fast and Lightweight Resemblance Detection Approach for Secondary Storage" (IEEE, 2018)}
\textbf{Ссылка:} \url{https://ieeexplore.ieee.org/document/8416390} \\
\textbf{Нейросетевой обзор:} Работа описывает алгоритм Odess, предназначенный для быстрого обнаружения сходства данных при дельта-компрессии. Метод использует контентно-зависимую выборку (Content-Defined Sampling, CDS) для формирования небольшого "эскиза" (sketch) данных. Это кратно ускоряет поиск похожего базового чанка по сравнению с традиционными методами, такими как N-Transform, сохраняя при этом высокий коэффициент сжатия. \\
\textbf{Зачем это нужно в работе:} В проекте исследуется и применяется \texttt{OdessHasher}. Статья необходима для понимания математического аппарата алгоритма и поиска путей его алгоритмической оптимизации при реализации.

\subsection*{3. "Finesse: Fine-Grained Feature Locality based Fast Resemblance Detection for Post-Deduplication Delta Compression"}
\textbf{Ссылка:} \url{https://www.usenix.org/conference/fast19/presentation/zhang} \\
\textbf{Нейросетевой обзор:} Авторы представляют подход Finesse, который группирует признаки (features) блоков данных в "супер-признаки" (super-features) для ускоренного поиска похожих блоков. Использование локальности признаков позволяет системе эффективно фильтровать кандидатов для дельта-компрессии с минимальными затратами оперативной памяти. \\
\textbf{Зачем это нужно в работе:} Идеи группировки и фильтрации из Finesse могут послужить источником вдохновения для создания новых эвристик в \texttt{marline}, позволяющих отсекать заведомо плохие кандидаты на этапе поиска базового чанка.

\subsection*{4. Data Deduplication and Delta Compression Concepts}
\textbf{Ссылка:} \url{https://www.ssrc.ucsc.edu/Papers/muthitacharoen-sosp01.pdf} (A Low-Bandwidth Network File System) \\
\textbf{Нейросетевой обзор:} Фундаментальная работа по алгоритму LBFS, впервые популяризировавшая Content-Defined Chunking (CDC). Рассматриваются базовые принципы разбиения данных на блоки на основе их содержимого. \\
\textbf{Зачем это нужно в работе:} Полезно для написания вводной части отчёта, так как SBC системы базируются на CDC. Статья помогает описать разницу между поиском точных дубликатов (CDC) и сходств (SBC).

\end{document}
